% Compiler avec pdfLaTeX (deux passages pour la table des matières).
% Les chemins d'images sont valides depuis le dossier livrables_jury/05_source_latex/.
\documentclass[11pt,a4paper]{article}

\usepackage[T1]{fontenc}
\usepackage[utf8]{inputenc}
\usepackage[french]{babel}
\usepackage{lmodern}
\usepackage[a4paper,margin=2.2cm]{geometry}
\usepackage{graphicx}
\usepackage{booktabs}
\usepackage{longtable}
\usepackage{tabularx}
\usepackage{array}
\usepackage{float}
\usepackage{xcolor}
\usepackage{enumitem}
\usepackage{fancyhdr}
\usepackage{hyperref}
\usepackage{titlesec}

\definecolor{atlasgreen}{HTML}{08783F}
\definecolor{atlaslight}{HTML}{E8F5EC}
\definecolor{atlasorange}{HTML}{A86406}
\definecolor{atlasred}{HTML}{B53C36}
\hypersetup{colorlinks=true,linkcolor=atlasgreen,urlcolor=atlasgreen}
\graphicspath{{../../assets/captures/preuves/}{../../assets/captures/fausses_pistes/}{../../assets/captures/bruits/}{../../assets/captures/decisions/}{../../assets/captures/cryptographie/}}
\setlength{\parskip}{0.45em}
\setlength{\parindent}{0pt}
\renewcommand{\arraystretch}{1.2}
\titleformat{\section}{\Large\bfseries\color{atlasgreen}}{\thesection}{0.6em}{}
\titleformat{\subsection}{\large\bfseries\color{atlasgreen}}{\thesubsection}{0.5em}{}
\pagestyle{fancy}
\fancyhf{}
\lhead{AtlasGrid -- MIRAGE}
\rhead{Exercice KASBAH -- document fictif}
\cfoot{\thepage}

\newcommand{\capture}[2]{%
  \begin{figure}[H]
    \centering
    \includegraphics[width=0.94\textwidth]{#1}
    \caption{#2}
  \end{figure}
}
\newcommand{\preuve}{\textcolor{atlasgreen}{\textbf{Preuve}}}
\newcommand{\bruit}{\textcolor{atlasorange}{\textbf{Bruit}}}
\newcommand{\faussepiste}{\textcolor{atlasred}{\textbf{Fausse piste}}}
\newcommand{\decision}{\textcolor{atlasgreen}{\textbf{Décision}}}

\begin{document}

\begin{center}
  {\Huge\bfseries\color{atlasgreen} Rapport d'incident AtlasGrid -- MIRAGE\par}
  \vspace{0.4cm}
  {\Large Chronologie, preuves, bruits, fausses pistes et décisions\par}
  \vspace{0.3cm}
  \textit{Exercice KASBAH -- entreprise fictive -- version consolidée Actes I, II et III}
\end{center}

\vspace{0.4cm}
\colorbox{atlaslight}{\parbox{0.94\textwidth}{
\textbf{Synthèse.} MIRAGE a chiffré 23 serveurs sur 40. Les journaux établissent 117,8 Go de transferts sortants vers une infrastructure externe. Une publication SIROCCO corrobore l'exposition de données clients et contractuelles ; la revendication de 300 Go demeure non confirmée. La cellule refuse de payer et reconstruit depuis l'air-gap de Settat, avec une remise en service progressive.}}

\tableofcontents
\clearpage

\section{Règles de lecture et périmètre}

\begin{itemize}[leftmargin=*]
  \item Une \preuve{} est un fait établi par une pièce conservée. Certaines pièces externes peuvent rester \emph{à confirmer} sur leur périmètre exact.
  \item Une \faussepiste{} est un signal crédible au départ, vérifié puis écarté par des éléments factuels.
  \item Un \bruit{} est un événement réel, mais attendu ou insuffisamment corrélé à MIRAGE.
  \item Une \decision{} est un arbitrage traçable : elle indique le risque traité, l'alternative écartée et le résultat attendu.
\end{itemize}

Le rapport ne transforme pas une hypothèse en certitude. En particulier, le phishing A-11 est réel mais n'est pas la porte d'entrée technique démontrée d'AtlasGrid ; l'attribution réelle de SIROCCO n'est pas établie au-delà de la revendication et de l'échantillon publié.

\section{Chronologie consolidée par jour}

\begin{longtable}{p{2.3cm}p{9.8cm}p{2.1cm}}
\toprule
\textbf{Moment} & \textbf{Événement établi} & \textbf{Pièces} \\
\midrule
\endhead
J-42 & Un rapport préliminaire OasisNet décrit un hameçonnage d'un technicien, l'accès à la console d'infogérance et l'extraction possible du secret partagé \texttt{svc\_oasisnet}. Le périmètre côté prestataire reste à confirmer. & A-08 \\
J-21 à J-1 & Des sessions VPN de \texttt{svc\_oasisnet} se produisent la nuit, sans MFA et depuis des IP incompatibles avec le profil attendu. & A-02, A-10 \\
J-11, 02:03 & Création de la règle pare-feu \texttt{OUT-TEMP-443} depuis FIN-112 vers HTTPS, avec journalisation désactivée. & A-09 \\
J-10 à J-1 & FIN-112 puis FILER-RBT transfèrent 117,8 Go vers \texttt{45.137.184.62:443}. & A-03 \\
J-3, 01:12 & La rétention est modifiée ; les travaux de sauvegarde J-3, J-2 et J-1 sont exclus. & A-05 \\
J-1, 03:11--03:16 & Privilèges spéciaux pour \texttt{svc\_oasisnet}, lancement de \texttt{svhost32.exe}, désactivation de Defender, C2, chiffrement des partages et suppression de VSS/catalogue. & A-01, A-04, A-09, A-10, A-12 \\
J1 & Phishing reçu, puis dépôts BKP-01/BKP-02 indisponibles. Le phishing est confirmé mais le lien causal avec MIRAGE ne l'est pas. & A-11, A-05 \\
J2 & 23/40 serveurs sont chiffrés ; une fraude BEC est bloquée ; SIROCCO publie un échantillon de données. & A-13, A-23, A-07 \\
J3 & Les sauvegardes en ligne récentes sont suspectes ou infectées. L'air-gap de Settat est retenu pour la reprise. & A-06, A-30 \\
J3, reprise & Refus de payer, restauration de l'ERP en priorité et réouverture client par paliers. & Décisions 9 à 12 \\
\bottomrule
\end{longtable}

\section{Preuves regroupées par phase de l'attaque}

\subsection{J-42 -- vecteur initial rapporté}
\subsubsection*{A-08 -- Rapport préliminaire d'OasisNet \preuve{} à confirmer}
Le CSIRT OasisNet indique qu'un technicien a saisi ses identifiants sur un portail de migration contrefait. Ces identifiants auraient donné accès à la console d'infogérance, d'où le secret de \texttt{svc\_oasisnet} a pu être extrait. Cette explication recoupe les journaux internes sur le compte utilisé, mais le rapport demeure préliminaire : son étendue et les responsabilités exactes restent ouvertes.
\capture{36_A08_rapport_preliminaire_OasisNet.png}{A-08 -- rapport préliminaire du CSIRT OasisNet.}

\subsection{J-21 à J-1 -- accès, préparation et exfiltration}
\subsubsection*{A-02 -- Journal VPN \preuve{}}
Le compte prestataire \texttt{svc\_oasisnet} ouvre plus de 90 sessions RemoteInteractive entre environ 02:00 et 05:00, sans MFA et depuis des IP hors profil. Ces comportements ne correspondent ni aux horaires ni au mode normal d'administration OasisNet.
\capture{04_A02_VPN_svc_oasisnet_anormal.png}{A-02 -- accès VPN anormaux de \texttt{svc\_oasisnet}.}

\subsubsection*{A-10 -- Journal d'authentification SIEM \preuve{}}
Le SIEM confirme les ouvertures RemoteInteractive, les privilèges spéciaux attribués à J-1 03:11, la création de \texttt{svhost32.exe} et l'effacement ultérieur du journal de sécurité. Il relie le compte de service à la phase d'exécution, au-delà des seules sessions VPN.
\capture{13_A10_SIEM_acces_svc_oasisnet.png}{A-10 -- traces SIEM corrélées au compte de service.}

\subsubsection*{A-09 -- Pare-feu et canal C2 \preuve{}}
À J-11 02:03, \texttt{svc\_oasisnet} crée \texttt{OUT-TEMP-443}, une règle sortante sans journalisation détaillée. Pendant l'exécution, FIN-112 contacte régulièrement \texttt{45.137.184.62:443} ; la préparation réseau et la communication de commande sont cohérentes.
\capture{10_A09_regle_firewall_C2.png}{A-09 -- règle pare-feu sortante et canal C2.}

\subsubsection*{A-03 -- Journal proxy d'exfiltration \preuve{}}
Sur dix nuits, FIN-112 puis FILER-RBT transfèrent 117,8 Go vers la même destination TLS, avec SNI et empreinte JA3 cohérents. Le volume est objectivé ; le contenu de chaque flux et les 300 Go annoncés par SIROCCO ne le sont pas par cette seule pièce.
\capture{14_A03_proxy_exfiltration_117_8Go.png}{A-03 -- exfiltration objectivée à 117,8 Go.}

\subsection{J-3 -- sabotage des sauvegardes}
\subsubsection*{A-05 -- Journal des travaux de sauvegarde \preuve{}}
La politique de rétention est modifiée par \texttt{svc\_oasisnet} à J-3 01:12. Les travaux les plus récents sont exclus, puis les dépôts BKP-01 et BKP-02 deviennent injoignables. Cette action réduit délibérément les possibilités de restauration avant le chiffrement.
\capture{08_A05_sauvegardes_modifiees_indisponibles.png}{A-05 -- sabotage de la rétention et indisponibilité des dépôts.}

\subsection{J-1 -- exécution et chiffrement}
\subsubsection*{A-01 -- Alerte EDR sur FIN-112 \preuve{}}
À 03:12:09, Defender détecte \texttt{Ransom:Win32/Mirage.A} dans le binaire non signé \texttt{C:\textbackslash{}Windows\textbackslash{}svhost32.exe}. La quarantaine échoue ; l'EDR observe ensuite une écriture de masse et l'altération de la protection temps réel. FIN-112 est donc le patient zéro établi.
\capture{02_A01_EDR_FIN-112_MIRAGE.png}{A-01 -- détection MIRAGE sur FIN-112.}

\subsubsection*{A-12 -- Chronologie EDR de FIN-112 \preuve{}}
La séquence EDR précise les actions : désactivation de Defender, lecture de 9\,412 fichiers, communication C2, renommage en \texttt{.mirage}, suppression des clichés VSS, modification de BCD et suppression du catalogue de sauvegarde. Ces traces reconstituent le mode opératoire sans dépendre d'une revendication extérieure.
\capture{15_A12_chronologie_EDR_FIN-112.png}{A-12 -- chronologie technique de l'exécution et du chiffrement.}

\subsubsection*{A-04 -- Fiche Forensic \texttt{LISEZMOI\_MIRAGE.txt} \preuve{}}
La note de rançon est créée sur FILER-RBT-02 autour de 03:15. Les attributs NTFS présentent des dates incohérentes avec les traces de création et d'accès, ce qui impose la conservation des métadonnées et interdit de s'appuyer sur les seules dates affichées.
\capture{11_A04_metadonnees_MIRAGE_alterees.png}{A-04 -- métadonnées NTFS de la note MIRAGE.}

\subsection{J1 et J2 -- impact, exposition et incident parallèle}
\subsubsection*{A-11 -- Courriel de hameçonnage \preuve{}, causalité non établie}
Le message reçu de \texttt{atlasgrid-it.info} échoue SPF et DMARC, ne présente pas de DKIM et invite à saisir des identifiants. Il doit être conservé et traité, mais ne constitue pas à lui seul la porte d'entrée démontrée de MIRAGE.
\capture{01_A11_entetes_phishing_atlasgrid-it.png}{A-11 -- en-têtes du phishing AtlasGrid.}

\subsubsection*{A-13 -- Inventaire d'impact \preuve{}}
L'inventaire SCCM confirme 23 serveurs chiffrés sur 40, dont PAIE-01, FACT-02, ERP-APP-01, ERP-DB-01 et les serveurs de fichiers. OT/SCADA demeure isolé et intact : il faut donc protéger cette situation plutôt que provoquer une coupure inutile.
\capture{12_A13_inventaire_serveurs_impact.png}{A-13 -- périmètre de l'impact confirmé.}

\subsubsection*{A-23 -- Fraude au président (BEC) \preuve{}, incident distinct}
Une usurpation de la PDG demande un virement de 480\,000 MAD vers un IBAN étranger. Le virement est bloqué. L'événement est réel et doit être préservé, mais aucun élément ne l'intègre à la chaîne technique MIRAGE.
\capture{18_A23_BEC_fraude_au_president.png}{A-23 -- tentative BEC bloquée.}

\subsubsection*{A-07 -- Publication SIROCCO \preuve{} d'exposition}
Le site SIROCCO publie un échantillon de 2\,400 lignes de données clients et contractuelles. Il corrobore l'exposition observée dans les journaux proxy. Les 300 Go affichés restent une revendication non confirmée et ne doivent pas être annoncés comme un fait.
\capture{31_A07_leaksite_SIROCCO_donnees_clients.png}{A-07 -- échantillon publié par SIROCCO.}

\subsection{J3 -- choix de la source de reprise}
\subsubsection*{A-06 -- Inventaire des sauvegardes \preuve{}}
Le registre DSI indique que la sauvegarde en ligne est hors service depuis J-3. La copie trimestrielle LTO-9 est stockée à Settat, hors réseau, dans un coffre ignifugé ; elle date de J-42 et n'a jamais été testée. Cette pièce établit une source de reprise indépendante, mais pas une réussite automatique de restauration.
\capture{32_A06_registre_moyens_sauvegarde.png}{A-06 -- registre des moyens de sauvegarde.}
\capture{34_A06_copie_airgap_Settat.png}{A-06 -- détail de la copie air-gap de Settat.}

\subsubsection*{A-30 -- Intégrité des sauvegardes en ligne \preuve{}}
Veeam indique RP-J-2 sain, RP-J-1 suspect et RP-J-0 infecté. Le même chargeur \texttt{svhost32.exe} que dans A-01 est présent dans les derniers points en ligne. Restaurer rapidement depuis ces points risquerait de réintroduire MIRAGE ; l'air-gap de Settat n'est pas concerné.
\capture{35_A30_integrite_sauvegardes_en_ligne.png}{A-30 -- analyse d'intégrité Veeam.}

\clearpage
\section{Fausses pistes écartées}

\subsection*{A-19 -- Clé USB du parking \faussepiste{}}
La clé USB trouvée sur le parking est analysée en sandbox mais n'a jamais été connectée au SI AtlasGrid. Elle n'apporte aucun lien exploitable avec l'attaque.
\capture{03_A19_cle_USB_parking_sandbox.png}{A-19 -- clé USB analysée et écartée.}

\subsection*{A-20 -- Courriel de migration \faussepiste{}}
Le message de migration interne est authentique : expéditeur, SPF/DKIM/DMARC et intranet sont cohérents. Il ne demande pas d'identifiants ; sa ressemblance avec un phishing ne suffit pas à le retenir.
\capture{05_A20_courriel_migration_legitime.png}{A-20 -- courriel de migration légitime.}

\subsection*{A-14 -- Badge datacenter \faussepiste{}}
Les horodatages de badge étaient fragiles parce que le lecteur était en maintenance. Sans corrélation AD, EDR ou caméra exploitable, le badge ne démontre aucune intrusion physique.
\capture{06_A14_badge_DATACENTER_non_fiable.png}{A-14 -- badgeuse non fiable.}

\subsection*{A-17 -- Copie USB de paie \faussepiste{}}
La copie locale de paie est autorisée par un ticket, réalisée en heures ouvrées et sans sortie réseau. Le contexte administratif documenté écarte une exfiltration.
\capture{07_A17_USB_paie_sauvegarde_autorisee.png}{A-17 -- opération de paie autorisée.}

\subsection*{A-16 -- Pic de trafic après presse \faussepiste{}}
Le pic correspond à la publication d'un article de presse, avec visiteurs humains et cache efficace. Il ne révèle pas un DDoS ni une phase technique de MIRAGE.
\capture{17_A16_pic_trafic_presse_legitime.png}{A-16 -- trafic web légitime après couverture presse.}

\subsection*{A-15 -- Jeton d'ex-salarié \faussepiste{}}
Le jeton mobile résiduel d'un ex-salarié est un défaut d'hygiène à corriger, mais les journaux ne montrent ni action privilégiée ni activité dans la fenêtre d'attaque.
\capture{19_A15_jeton_ex_salarie_sans_lien.png}{A-15 -- jeton résiduel sans lien démontré.}

\subsection*{A-24 -- Salarié en litige \faussepiste{}}
Les sources VPN, AD et badge ainsi que l'arrêt de travail documenté excluent le salarié de la fenêtre d'attaque. L'hypothèse interne doit donc être abandonnée.
\capture{26_A24_salarie_litige_hors_de_cause.png}{A-24 -- piste interne écartée.}

\subsection*{A-18 -- Revendication DarkAtlas \faussepiste{}}
DarkAtlas revendique l'incident, mais ses données sont publiques et ses indicateurs ne correspondent pas à MIRAGE : extension, chronologie et canal de rançon divergent. Une revendication externe n'est jamais une attribution suffisante.
\capture{27_A18_revendication_DarkAtlas_extrait.png}{A-18 -- revendication initiale à investiguer.}
\capture{28_A18_revendication_DarkAtlas_invalidee.png}{A-18 -- revendication invalidée après recoupement.}

\clearpage
\section{Bruits écartés}

\subsection*{A-21 -- Scans Internet génériques \bruit{}}
Les scans proviennent de moteurs de mesure Internet et ne réalisent que des SYN ; aucun actif interne n'est atteint. L'événement est réel, mais non corrélé à MIRAGE.
\capture{20_A21_scans_internet_generiques.png}{A-21 -- scans génériques sans compromission.}

\subsection*{A-25 -- PsExec de maintenance \bruit{}}
PsExec est officiel, signé et associé au ticket de maintenance \#4502. La procédure attendue ne correspond pas à une exécution malveillante.
\capture{21_A25_PsExec_maintenance_autorisee.png}{A-25 -- opération PsExec autorisée.}

\subsection*{A-27 -- Archivage trimestriel \bruit{}}
L'archivage est prévu par le ticket \#4381 ; les fichiers restent relisibles et non chiffrés. Il ne doit pas être confondu avec le chiffrement MIRAGE.
\capture{22_A27_archivage_trimestriel_autorise.png}{A-27 -- archivage autorisé.}

\subsection*{A-26 -- Connexion depuis Paris \bruit{}}
La connexion est réalisée depuis l'appareil connu du directeur commercial, dans un contexte de déplacement et d'activité de messagerie normale.
\capture{24_A26_connexion_Paris_legitime.png}{A-26 -- connexion légitime.}

\subsection*{A-40 -- Ticket OasisNet d'un autre client \bruit{}}
Le ticket mentionne une extension de rançongiciel, mais concerne un autre client et aucun actif AtlasGrid. Il n'a aucune valeur d'attribution pour l'incident.
\capture{25_A40_ticket_OasisNet_autre_client.png}{A-40 -- ticket sans lien avec AtlasGrid.}

\subsection*{A-41 -- Assistance à distance \bruit{}}
La session est autorisée, l'opérateur identifié, le ticket HELP-3391 présent et le MFA valide. Elle est distincte des sessions anormales observées sur \texttt{svc\_oasisnet}.
\capture{29_A41_assistance_distance_autorisee.png}{A-41 -- assistance distante autorisée.}

\subsection*{A-22 -- Faux positif CoinMiner \bruit{}}
L'alerte vise 7-Zip signé pendant une compilation nocturne planifiée. L'analyse ne montre ni minage ni lien avec MIRAGE.
\capture{30_A22_faux_positif_cryptominer.png}{A-22 -- faux positif CoinMiner.}

\subsection*{A-28 -- Compte \texttt{svc\_backup} \bruit{}}
Le compte réalise chaque nuit une tâche Veeam planifiée depuis VBR-01, avec un profil stable depuis des mois. Ses paramètres diffèrent de ceux de \texttt{svc\_oasisnet}.
\capture{31_A28_svc_backup_planifie.png}{A-28 -- tâche de sauvegarde attendue.}

\subsection*{A-31 -- Alerte macro DG-01 \bruit{}}
L'alerte date de J-90, concerne un fichier de 2019 déjà mis en quarantaine et ne présente aucune sortie réseau associée. Elle ne rejoint pas le C2 nocturne.
\capture{32_A31_macro_DG01_historique.png}{A-31 -- alerte macro historique, déjà traitée.}

\subsection*{A-29 -- Trafic Microsoft \bruit{}}
Les domaines Microsoft sont contactés en heures ouvrées depuis l'ensemble des postes pour Windows Update et la télémétrie. Les destinations et horaires diffèrent du canal \texttt{45.137.x.x}.
\capture{33_A29_trafic_Microsoft_legitime.png}{A-29 -- trafic Microsoft légitime.}

\clearpage
\section{Décisions prises et justification}

\subsection{Acte I -- confinement}
\subsubsection*{Décision 1 -- Isoler FIN-112 sous tension}
\decision{} : isoler FIN-112 du réseau sans l'éteindre. Cela stoppe la propagation tout en conservant la mémoire, les processus et les connexions utiles à l'investigation. Éteindre immédiatement aurait fait disparaître des traces volatiles ; ne rien isoler aurait laissé le chiffrement se propager.
\capture{01_isoler_FIN-112_sous_tension.png}{Décision 1 -- isolement de FIN-112 sous tension.}

\subsubsection*{Décision 2 -- Isoler les segments critiques}
\decision{} : segmenter les flux à risque et conserver sous contrôle les services indispensables. Une coupure générale aurait été disproportionnée ; l'absence de segmentation aurait facilité les mouvements latéraux SMB déjà observés.
\capture{02_isoler_segments_critiques.png}{Décision 2 -- confinement ciblé des segments critiques.}

\subsection{Acte II -- communication, continuité et conformité}
\subsubsection*{Décision 3 -- Communiquer en interne sans détails sensibles}
\decision{} : donner des consignes aux salariés et éviter le silence, tout en ne révélant ni indicateurs techniques ni données personnelles. Le silence nourrit les rumeurs ; une communication non vérifiée peut compromettre l'enquête.
\capture{03_communication_interne_rassurer_sans_detail.png}{Décision 3 -- communication interne factuelle.}

\subsubsection*{Décision 4 -- Maintenir l'OT sous surveillance renforcée}
\decision{} : l'OT/SCADA reste actif, isolé et surveillé. L'absence de session compromise ne justifie pas une coupure préventive à fort impact industriel, mais des critères de bascule sont préparés.
\capture{04_OT_surveillance_renforcee_sans_coupure.png}{Décision 4 -- surveillance OT sans coupure immédiate.}

\subsubsection*{Décision 5 -- Déclarer l'assureur immédiatement}
\decision{} : déclarer dans le délai de 48 heures de l'exercice, préserver les preuves et demander les validations nécessaires aux dépenses majeures. Attendre risquerait d'affecter la couverture.
\capture{05_declaration_assureur_immediate.png}{Décision 5 -- déclaration immédiate à l'assureur.}

\subsubsection*{Décision 6 -- Répondre à la presse avec une version vérifiée}
\decision{} : répondre de manière factuelle à Maghreb Éco, sans confirmer les 300 Go ni divulguer de données ou de détails d'enquête. Le résultat est un article équilibré : répondre a évité que la revendication de l'attaquant devienne le seul récit public.
\capture{06_communication_presse_repondre_version.png}{Décision 6 -- réponse presse vérifiée.}
\capture{14_retour_presse_et_validation_airgap.png}{Retour presse : l'article est équilibré et la stratégie de reprise est validée.}

\subsubsection*{Décision 7 -- Répondre factuellement au client}
\decision{} : informer Chérifienne des Mines des mesures et du calendrier de mise à jour sans attester une absence d'exposition non prouvée. Une promesse absolue aurait créé un risque contractuel et réputationnel.
\capture{07_reponse_client_rester_factuel.png}{Décision 7 -- réponse client factuelle.}

\subsubsection*{Décision 8 -- Notifier l'Autorité}
\decision{} : notifier avec les faits connus et les incertitudes, puis compléter l'analyse. La démarche respecte le délai de 72 heures propre au scénario et protège les personnes potentiellement concernées.
\capture{08_notification_autorite_transparente.png}{Décision 8 -- notification transparente à l'Autorité.}

\subsection{Acte III -- restauration et retour des clients}
\subsubsection*{Décision 9 -- Ne pas payer la rançon}
\decision{} : ne pas financer l'extorsion car aucun déchiffrement ni effacement de données n'est garanti. A-06 et A-30 offrent une alternative de reprise indépendante. Le paiement et la restauration en ligne rapide sont donc écartés.
\capture{09_direction_ne_pas_payer.png}{Décision 9 -- refus de paiement.}

\subsubsection*{Décision 10 -- Restaurer le cœur ERP en priorité}
\decision{} : restaurer la base et l'applicatif ERP avant la paie et la facturation, qui dépendent de ce socle. Le premier résultat visible est plus tardif, mais le redémarrage est plus complet et cohérent.
\capture{11_priorite_restauration_ERP.png}{Décision 10 -- priorité au cœur ERP.}
\capture{13_sequence_reprise_ERP.png}{Séquence opérationnelle de reprise ERP.}

\subsubsection*{Décision 11 -- Restaurer depuis l'air-gap de Settat}
\decision{} : utiliser les bandes LTO-9 déconnectées, malgré leur ancienneté et le délai estimé. Les points en ligne récents sont suspect ou infecté ; l'air-gap réduit le risque de réintroduire MIRAGE.
\capture{12_choix_airgap_Settat.png}{Décision 11 -- choix de l'air-gap de Settat.}
\capture{10_plan_reprise_airgap_Settat.png}{Plan de restauration depuis Settat : AD, ERP, paie/facturation, partages.}

\subsubsection*{Décision 12 -- Rouvrir les services clients par paliers}
\decision{} : reconnecter progressivement après validation de sécurité, d'intégrité et de stabilité à chaque étape. Cette option évite à la fois une promesse de reprise complète prématurée et un blocage inutile des services prêts à être ouverts.
\capture{15_reprise_progressive_client_cherifienne.png}{Décision 12 -- reprise client progressive et encadrée.}

\section{Cadre de conformité et conclusion}

Les décisions sont fondées sur les faits conservés et sur une logique de proportionnalité : préserver les preuves, contenir, communiquer sans spéculer, restaurer dans un environnement assaini et documenter chaque arbitrage. Cette méthode s'inscrit dans l'esprit du NIST CSF 2.0 et du NIST SP 800-61 Rev. 3. Dans le cadre pédagogique AtlasGrid, elle prend aussi en compte la loi marocaine n\textsuperscript{o} 09-08 pour la protection des données et la logique de sécurité des systèmes sensibles évoquée par la loi n\textsuperscript{o} 05-20.

\textbf{Conclusion à présenter au jury.} La chaîne technique la plus étayée est : compromission possible du prestataire OasisNet $\rightarrow$ usage anormal de \texttt{svc\_oasisnet} $\rightarrow$ règle de sortie $\rightarrow$ exfiltration et sabotage des sauvegardes $\rightarrow$ exécution sur FIN-112 $\rightarrow$ chiffrement MIRAGE $\rightarrow$ impact métier et exposition corroborée. La reprise est volontairement plus lente, mais indépendante de l'infrastructure compromise, contrôlée et traçable.

\end{document}
